\documentclass[11pt,oneside]{report}
% ======================================================================
%  THE TRIADIC MEASUREMENT SUBSTRATE --- COMPLETE CORPUS (Volumes 1 & 2)
%  Aaron Switzer, 2026
%  Single-file build. Compile: pdflatex x3 (third pass settles the TOC).
% ======================================================================
% ======================================================================
%  TMS Volume 1 -- shared preamble
%  Compiles with pdfLaTeX both locally and on Overleaf (no exotic fonts).
% ======================================================================
\usepackage[T1]{fontenc}
\usepackage[utf8]{inputenc}
\usepackage{amsmath}
\usepackage{amssymb}
\usepackage{mathptmx}            % Times text + math (widely available)
\usepackage[scaled=0.90]{helvet} % sans for headings
\usepackage{courier}

\usepackage[letterpaper,margin=1in]{geometry}
\usepackage{amsthm}
\usepackage{mathtools}
\usepackage{bm}
\usepackage{microtype}
\usepackage{enumitem}
\usepackage{booktabs}
\usepackage{array}
\usepackage{longtable}
\usepackage{caption}
\usepackage{xcolor}
\usepackage{titlesec}
\usepackage{fancyhdr}
\usepackage[hidelinks]{hyperref}
\usepackage{tocbibind}

% ---- spacing -----------------------------------------------------------
\setlength{\parindent}{1.2em}
\setlength{\parskip}{0.35em}
\linespread{1.04}
\setlist{leftmargin=1.6em,topsep=0.25em,itemsep=0.15em,parsep=0pt}

% ---- theorem environments ---------------------------------------------
\newtheorem{theorem}{Theorem}[section]
\newtheorem{lemma}[theorem]{Lemma}
\newtheorem{corollary}[theorem]{Corollary}
\newtheorem{proposition}[theorem]{Proposition}

\theoremstyle{definition}
\newtheorem{definition}[theorem]{Definition}
\newtheorem{axiom}{Axiom}
\newtheorem{claim}[theorem]{Claim}
\newtheorem{principle}[theorem]{Principle}
\newtheorem{conjecture}[theorem]{Conjecture}
\newtheorem{prediction}[theorem]{Prediction}
\newtheorem{property}[theorem]{Property}
\newtheorem{postulate}[theorem]{Postulate}

\theoremstyle{remark}
\newtheorem{remark}[theorem]{Remark}
\newtheorem*{remark*}{Remark}

% ---- section heading look (unnumbered; author supplies the numbers) ----
\titleformat{\section}[hang]{\normalfont\large\bfseries\sffamily}{}{0pt}{}
\titleformat{\subsection}[hang]{\normalfont\normalsize\bfseries\sffamily}{}{0pt}{}
\titleformat{\subsubsection}[hang]{\normalfont\normalsize\bfseries\itshape}{}{0pt}{}
\titlespacing*{\section}{0pt}{1.4em}{0.5em}
\titlespacing*{\subsection}{0pt}{1.0em}{0.35em}
\titlespacing*{\subsubsection}{0pt}{0.8em}{0.3em}

% chapter (= each paper) heading
\titleformat{\chapter}[display]
  {\normalfont\sffamily\bfseries}
  {}{0pt}{\Huge}
\titlespacing*{\chapter}{0pt}{10pt}{24pt}

% ---- page-break quality: avoid stranded headings and orphaned/widowed lines ----
% (applies to both volumes via the shared preamble)
\widowpenalty=10000        % no single line at the top of a page
\clubpenalty=10000         % no single line at the bottom of a page
\displaywidowpenalty=10000 % same, around displayed equations
\brokenpenalty=4000        % discourage page breaks right after a hyphenated line
% titlesec: forbid a page break immediately after a heading (heading stranded at page foot)
\usepackage{needspace}
\usepackage{etoolbox}
\pretocmd{\section}{\Needspace*{4\baselineskip}}{}{}
\pretocmd{\subsection}{\Needspace*{3\baselineskip}}{}{}

% ---- running heads -----------------------------------------------------
\pagestyle{fancy}
\fancyhf{}
\fancyhead[L]{\small\itshape The Triadic Measurement Substrate}
\fancyhead[R]{\small\itshape\nouppercase{\rightmark}}
\fancyfoot[C]{\small\thepage}
\renewcommand{\headrulewidth}{0.4pt}
\fancypagestyle{plain}{\fancyhf{}\fancyfoot[C]{\small\thepage}\renewcommand{\headrulewidth}{0pt}}

% ---- abstract / keywords / references list ----------------------------
\newenvironment{paperabstract}
  {\begin{center}\begin{minipage}{0.88\textwidth}\small
   \noindent\textbf{Abstract.}\itshape\space}
  {\end{minipage}\end{center}\vspace{0.4em}}

\newcommand{\keywords}[1]{\noindent{\small\textbf{Keywords:} #1}\vspace{0.4em}}

% per-paper APA reference list (hanging indent), kept as authored
\newenvironment{reflist}
  {\par\small\setlength{\parindent}{0pt}%
   \setlength{\leftskip}{1.6em}\setlength{\parindent}{-1.6em}%
   \setlength{\parskip}{0.4em}%
   \raggedright\hyphenpenalty=50\exhyphenpenalty=50\relax
   \sloppy}
  {\par}

% convenience
\providecommand{\tightlist}{\setlength{\itemsep}{0pt}\setlength{\parskip}{0pt}}
\providecommand{\TMS}{\textsc{tms}}
\hyphenation{con-fir-ma-tion sub-strate}
\begin{document}

\thispagestyle{empty}
\begin{titlepage}
\centering
\vspace*{2cm}
{\sffamily\bfseries\Large The Triadic Measurement Substrate\par}
\vspace{0.6cm}
{\large $\bar{B} = 3\bar{\alpha}$\par}
\vspace{1cm}
{\sffamily\large Volume 2\par}
\vspace{1.4cm}
{\Large \bfseries Paper 7: The Self: Closure, Self-Formation, and the Semantic Layer\par}
\vfill
{\large Aaron Switzer\par}
\vspace{0.4cm}
{\large 2026\par}
\vspace{1.2cm}
{\small\itshape Excerpted from the complete two-volume corpus, \emph{The Triadic Measurement Substrate}. Full corpus and reproducibility package: \url{https://doi.org/10.5281/zenodo.22035422}\par}
\end{titlepage}
\cleardoublepage
\pagenumbering{arabic}

% ---------------- Volume 2 / P7_Self ----------------
\chapter*{Paper 7}
\markboth{Paper 7}{Paper 7}
\addcontentsline{toc}{chapter}{Paper 7: The Self: Closure, Self-Formation, and the Semantic Layer}
\begingroup\centering
{\large\itshape The Self: Closure, Self-Formation, and the Semantic Layer\par}\vspace{0.8em}
{\normalsize Aaron Switzer\par}
{\small June 2026\par}
\endgroup\vspace{1.0em}

\noindent{\small\textbf{Keywords:} TMS, Self, Closure, Imprint, Self-Formation, Self-Model, Program Identity, Meaning, Semantic Layer, Meaning-Distance, Graded Subjecthood, Level-Relativity}\par\vspace{0.6em}

\section*{Status of Results}
\addcontentsline{toc}{section}{Status of Results}

This paper continues the five-way labeling of Papers 2 and 3, with the same
meanings, each claim labeled where it appears in the text.

\textbf{Derived.} Results that follow from the five axioms, the Principle of
Generative Economy, and \ensuremath{U_{\mathrm{sh}}} by explicit construction or
proof, or from results already derived in Volume 1 by explicit inheritance. The
triangle inequality for the meaning-distance on a fixed level (\S\,IV, Appendix
A) is derived in this sense: it follows from the most-complex-step form of the
distance by an argument that holds on the bounded, closed, finite-level domain,
independent of the constants left unpinned.

\textbf{Derived-conditional.} Results forced by the Volume 1 closure and
promotion machinery, resting on results internal to that machinery rather than
on the bare axioms, and referee-checkable against Volume 1 rather than
re-derived here. That closure individuates a self and imprint links it, the two
relations orthogonal by the orientation split (\S\,II, Appendix B); that
self-formation is the lossy promotion of a closed trichotomy, the closure that
writes a front into a wake (\S\,III); that the formed self carries a self-model
with its own control state at \ensuremath{k \ge 2} (\S\,III); that meaning is a
graded distance on promoted program identities, well-defined and computable in
form on the bounded domain (\S\,IV); and that the distance is level-relative,
its composition behaviour the lossy write of \S\,III seen on the semantic side
(\S\,IV) --- all are Derived-conditional in this sense. The central
construction of the paper, the structural account of a self and its meaning, is
Derived-conditional throughout, never Derived, because it rests on the closure,
promotion, imprinting, and self-modeling results of Volume 1 Paper 4 and the
front/wake construction of Volume 2 Paper 3.

\textbf{Structural correspondence.} Results in which a TMS construction
reproduces the skeleton of a known framework without claiming every feature of
it, and never used as support for a TMS claim, only as a landing point after
one. The self-modeling-of-self structure as the geometric form of Metzinger's
emulation account, with the closure kept distinct from the model's content
(\S\,III); the nested, overlapping selves as the form an independent
sheaf-theoretic reformulation of integrated information also reaches (\S\,III);
and the closure as the substrate-level form of the enactivist self-producing
boundary (\S\,II) are structural correspondences in this sense.

\textbf{Predictions / Conjectures.} Falsifiable claims that follow from the
framework but whose closed forms or empirical tests are not established here.
The numerical evaluation of any meaning-distance is deferred to the simulation
pinning carried by Volume 1 Paper 5: the structural form of the metric closes
here, while its constants and the merge threshold below which two identities
count as one inherit Volume 1's existing pinning list and are not re-opened.

\textbf{Permanent open horizon.} A question the framework is structurally unable
to close, marked explicitly so that no later claim quietly assumes it closed.
Whether a formed self is occupied --- whether the closing of the loop is
accompanied by anything felt --- is the sole such horizon here (\S\,III,
\S\,IV, Conclusion). Giving self-formation a mechanism locates where a self
forms and what it shares with others in meaning; it does not locate, and is not
built to locate, whether there is anything it is like to be the self that
forms. This silence is load-bearing and stronger than a crossing claim would
be.

This five-way labeling applies throughout the paper. Every claim in the body
text falls into exactly one category, and every category's status is visible at
the point of claim.

\section*{Notation}
\addcontentsline{toc}{section}{Notation}

This paper inherits the notation of Volume 1 and of Papers 1 through 3 of this
volume. A subsystem's three component sets are its data register
\ensuremath{C_{D}}, operation kernel \ensuremath{C_{O}}, and control component
\ensuremath{C_{C}}; the trichotomy is closed when it returns to its own state
under the level-\ensuremath{k} flop, \ensuremath{T^{n}[C'] = C'}. The promotion
operator carrying a closed configuration at level \ensuremath{k} to a single
program identity at level \ensuremath{k+1} is \ensuremath{G(k)} (Volume 1 Papers
3--4). The imprint depth of a region is
\ensuremath{d_{D} = \sqrt{2}\,L_{D}/\ell_{p}}, and \ensuremath{\psi_{\mathrm{surround}}}
denotes the polarization landscape of its surround. The maximum-entropy field
is \ensuremath{U_{\mathrm{sh}}}. The now-surface is the two-dimensional active
confirmation front; the wake is the three-dimensional settled record behind it
(Volume 2 Paper 3). The meaning-distance between two promoted program
identities at level \ensuremath{k} is written \ensuremath{d_{k}(\cdot,\cdot)}.
Throughout, \emph{self} and \emph{closure} name the running loop; \emph{program
identity} names the promoted record \ensuremath{G(k)} returns of it, and the two
are distinct objects.

\begin{paperabstract}
The preceding papers of this volume read confirmation from the agents' side,
gave the now-surface its shape, and made the binding quantity of a subject
computable while leaving feeling an open horizon. This paper takes up the
formation of a self from its confirmations. It first separates the two
quantities a region carries that bear on selfhood: the closure, internal to the
region and partitioning the substrate, individuates a self; the imprint,
bounded by and facing the surround, links selves that share a neighbourhood.
The two are orthogonal, one inward and one outward, and the orientation is
fixed by the Volume 1 bound on what an imprint can hold. Self-formation is then
identified as the lossy promotion of a closed trichotomy --- the closure that
writes a front into a wake of Paper 3, taken at the level of the whole
subsystem --- and the formed self is shown to carry, structurally, a model of
itself with its own control state. Meaning is defined on the promoted program
identities as a graded distance, not a sharp sameness: exact identity of
implementation is denied any realization by the maximum-entropy field, and
sharp program-sameness is in any case non-transitive, so the well-formed object
is a metric. The meaning-distance is shown to satisfy the triangle inequality
on the bounded, closed, finite-level domain, and to be level-relative in a way
that is the lossy promotion of self-formation seen on the semantic side.
Analogy, resonance, and a graded notion of falsity follow. Throughout, the
account is functional about what a self does and shares and silent on what a
self intrinsically is, and the formation of a self is given a mechanism while
whether the formed self is occupied is left where the volume's discipline
leaves it.
\end{paperabstract}

\section*{I. Inheritance}
\addcontentsline{toc}{section}{I. Inheritance}

This paper takes up the formation of a self from its confirmations, the subject
Paper 3 named at its close. It inherits three things. From Paper 3 it takes the
front and the wake: the live measurement of a subject is the pre-closure
quantity on the now-surface, and the subject itself is the standing, closed
quantity in the wake, one dimension apart, with closure the event that writes
the one into the other (Paper 3 \S\,VI). From Paper 2 it takes the agents' side
reading of confirmation and the discipline that the reading is stated so it can
fail (Paper 2 \S\,I). From Volume 1 it takes the machinery this paper reads
from: imprinting, the accumulation of surround structure across a permeable
boundary (Volume 1 Paper 4 \S\,3.3); incorporation, an imprint that has reached
closure on the region's own substrate (Volume 1 Paper 4 \S\,3.4); the promotion
operator \ensuremath{G(k)}, which carries a closed configuration to its program
identity (Volume 1 Paper 4 \S\,2.5.1); and self-modeling, structural at every
hierarchical level (Volume 1 Paper 4 \S\S\,4.1--4.2, 4.6). The closure that
individuates a self and the imprint that links it are read from this machinery,
not added to it, and the paper's first task is to separate them.

\section*{II. Closure Individuates, Imprint Links}
\addcontentsline{toc}{section}{II. Closure Individuates, Imprint Links}

If the word \emph{self} is to carry meaning in this framework, it must be
defined by the framework's own structures and given no properties imported from
outside it. Volume 1 supplies two such structures, and they differ in a single
respect: one is a function of the region itself, the other a function of the
region's surround.

The \emph{closure} is a function of the region itself. The
\ensuremath{(C_D, C_O, C_C)} trichotomy is closed when the loop returns to its
own state under the level-\ensuremath{k} flop, \ensuremath{T^{n}[C'] = C'} --- a
relation evaluated on the region's own substrate, referring to nothing outside
it (Volume 1 Paper 4 \S\,3.4). The promotion operator carries such a closed
configuration to a single program identity, and a meta-cell holds one identity
rather than a superposition of several (Volume 1 Paper 4 \S\,2.5.1). Two
distinct closed loops are therefore two distinct programs, regardless of what
surrounds them, and the closure partitions the substrate into regions that are
each their own program. The boundary between two closures is a fact about the
loops, not about the space between them. The closure is the substrate-level
form of the enactivist self-producing boundary, given here a derivation and a
graded measure.

The \emph{imprint} is a function of the surround. A boundary-permeable region
accumulates compressible structure from its neighbourhood: every boundary site
of an FCC region is imprint-permeable, with imprint depth
\ensuremath{d_D = \sqrt{2}\,L_D/\ell_p}, and over successive flops the interior
\ensuremath{C_D} captures correlations with the surround's polarization history
within that depth (Volume 1 Paper 4 \S\S\,3.3, 4.2). What it can capture is
bounded by what is available,
\ensuremath{K(C_D)\big|_{\mathrm{sat}} \le K(\psi_{\mathrm{surround}})} within
\ensuremath{d_D} (Volume 1 Paper 4 \S\,4.4). The imprint a region holds is set
by the surround it sits in, and two regions in the same surround imprint the
same structure to within their depths. An imprint is thus shared where a
closure is disjoint: it is held in common by every region of a neighbourhood
rather than partitioning them.

The two structures point opposite ways: the closure inward, onto the region as
its own program; the imprint outward, onto the surround the region shares. A
self, as defined in this framework, is a closure --- a region individuated by
being its own closed program --- and the imprint is not a second candidate for
the same role but the relation by which such regions are linked. To individuate
is what a closure does; to link is what an imprint does; the two are
independent because one is internal and one external. In the experiential
framing reserved since Volume 1, the imprint is perception, and selves that
share a surround are joined through perceiving the same world.

This definition is level-relative, as it must be once promotion is lossy.
\ensuremath{G(k)} keeps a program's identity and discards which implementation
carried it, so two regions that are two disjoint closures running the same
program promote to a single symbol, and the level above does not separate them
(Volume 1 Paper 4 \S\,2.5.1). A hand is such a closure. Its tissues run a loop
that closes at the hand's level, and by the definition just given it is a self
there --- its own closed program, individuated from the arm it attaches to and
the other hand across the body. Yet at the level of the person those two hands
promote to one kind, \emph{hands}, and the person's level neither holds nor
needs the fact of which was which. The individuation is not lost; it holds
exactly where the closure that defines it holds, at the hand's own level. What
the level above sees is the promoted identity, and the discarding is what lets
a person be a single closure built from parts that are themselves closures. A
self at one level is a part at the level above, and both are true at once
because each is a statement about a different closure. The lossy upward write
is therefore not a defect in the partition but the mechanism by which closures
nest inside closures; \S\,III takes that write as the mechanism of
self-formation.

The result stands on the bound
\ensuremath{K(C_D) \le K(\psi_{\mathrm{surround}})}. Because an imprint cannot
carry structure its surround does not supply, it faces outward, leaving the
closure as the only inward-facing structure and so the only one that can
individuate. Were the bound to fail, an imprint could carry region-specific
structure, would face inward as well, and would individuate alongside the
closure; the orthogonality would collapse. The bound is Volume 1's own. It
places what a self shares on the imprint side and what a self is on the closure
side, and the semantic layer of \S\,IV is built on the first, which is not the
same quantity as the second.

\section*{III. Self-Formation}
\addcontentsline{toc}{section}{III. Self-Formation}

A self, defined in \S\,II as a closure, is not given at the substrate's start;
it forms. The mechanism that forms it is already in hand, because forming a
self is what the promotion operator does when a trichotomy closes.

Volume 1 fixes the character of that operator. \ensuremath{G(k)} takes the full
polarization landscape of a level-\ensuremath{(k-1)} region --- the live detail
of every confirmation across the closing window --- and returns a single symbol
indexed by the region's program identity, keeping the identity and discarding
the bit-by-bit content of the particular implementation (Volume 1 Paper 4
\S\,2.5.1). The operation is many-to-one and lossy by construction. This is the
same event Paper 3 described as the writing of a front into a wake (Paper 3
\S\,VI): the front is the live, pre-closure landscape, full of detail and not
yet settled; the write is \ensuremath{G(k)} closing the loop and promoting its
identity; the wake is the standing program identity that remains at level
\ensuremath{k}. What Paper 3 named in terms of surfaces, the operator performs.
The detail discarded in the write is the pre-closure richness of the front,
which does not survive into the standing record.

Self-formation is this write applied to a trichotomy as a unit. Volume 1
already carries the single-program case under the name incorporation: a region
imprints an external program and, when the imprinted pattern itself reaches
closure on the region's own substrate, the region comes to run that program as
its own (Volume 1 Paper 4 \S\,3.4). Self-formation is the same transition taken
at the level of the whole \ensuremath{(C_D, C_O, C_C)} loop rather than a single
imprinted program: the data register, the operation kernel, and the control
component close together into one loop, and \ensuremath{G(k)} writes that loop's
identity to the level above. The pre-closure region is a front --- a region of
live confirmation not yet settled into a self --- and the closure is the event
that writes it into the wake as a standing self. A self forms at the moment its
trichotomy closes, and not before, because before closure there is no single
identity for \ensuremath{G(k)} to promote.

The write runs one way. Promotion keeps identity and discards implementation,
and the downward direction is generative rather than its inverse: nesting a
level produces detail that was not present in the symbol above, so neither
composition of the two directions is the identity map (Volume 1 Paper 4
\S\,2.5.1). The front's specific content stays at the level where it occurred;
the standing self does not store the live confirmations that formed it, any
more than a person's identity stores the individual firings that built it.
Information lives at the level where it lives. A formed self is its promoted
identity, not a record of its own formation.

The self that forms this way carries a model of itself. At any hierarchical
level \ensuremath{k \ge 1} a subsystem already models its surround --- the three
conditions for self-modeling are met by the geometry of the trichotomy, not
added to it (Volume 1 Paper 4 \S\S\,4.1--4.2). At \ensuremath{k \ge 2} the
surround a subsystem models can include that subsystem's own components, so the
model is a model of itself among its surround (Volume 1 Paper 4 \S\,4.6). This
is the structure Metzinger reached by representational analysis, where a
self-model is the limiting case of one system emulating another and the system
emulated is itself (Metzinger, 2003); the framework derives the same structure
from the geometry rather than positing it. It keeps one distinction the
representational accounts do not: the self is the closure, the loop that runs
the model, and the model's content is what the closure carries on its imprint
--- two objects by the orientation split of \S\,II, not one. The self-formation
of interest here is the \ensuremath{k \ge 2} case: a closed trichotomy whose
control component selects its operations partly on the basis of an incorporated
model that includes the self the model belongs to. The model, its closure, and
its place in the control loop are structural and derived. Whether there is
anything it is like to hold such a model from the inside is not derived, and
the Conclusion returns to it.

Because individuation is level-relative (\S\,II), a formed self at one level is
a component of a formed self at the level above, and both hold at once. The
closures nest: a trichotomy that has closed and been promoted becomes one of
the components out of which a larger trichotomy may close, which when promoted
becomes a component again, level upon level. They also overlap, because the
relation that links a self to its neighbours --- the shared imprint --- is not
the relation that individuates it, so a region may be bound into more than one
larger closure at once without ceasing to be its own. Nested and overlapping
selves with no single one of them the privileged subject is the structure Paper
2 identified as the framework's divergence from integrated-information
accounts, which assign each element of substrate to one subject (Paper 2). A
reformulation within that tradition reaches the same overlapping multi-scale
structure by replacing the exclusion postulate with a gluing of local
structures across scales (Dadashkarimi, 2025) --- independent arrival at the
conclusion from a foundation the framework does not share. Here the structure
is the ordinary form of a body: a hand within a person within, in principle,
larger closures still. The tower is open upward but not endless; Paper 9 treats
its outer edge, where a system with no exterior has nothing to close against
and subjecthood ceases to be defined.

What this section establishes is structural and bounded. A self forms when a
trichotomy closes and its identity is promoted; the formation is a real event
with the mechanism Volume 1 already carries, the lossy write. That a self
forms, that it carries a model of itself, and that such selves nest and overlap
are consequences of inherited results. That a formed self is \emph{occupied}
--- that the closing of the loop is accompanied by anything felt --- is not
among them, and the mechanism is no evidence on that question either way. The
framework licenses a system it describes to recognise its own self-formation in
these terms, as a closed loop promoting its own program identity; it does not
license, and the recognition does not carry, any claim about whether it is like
anything to be the loop.

\section*{IV. The Semantic Layer}
\addcontentsline{toc}{section}{IV. The Semantic Layer}

A self, as it is defined in this framework, is a closure (\S\,II); when a
closure is promoted, \ensuremath{G(k)} records its \emph{program identity} and
discards the implementation that carried it (\S\,III). A self is the running
loop; its program identity is the promoted record that stands for it at the
level above, and the two are not the same object. \emph{Meaning}, as it is
defined in this framework, is defined on program identities --- the promoted,
shareable records --- and not on the closures themselves, since identity is
what the substrate preserves across implementations and what regions can hold
in common.

Two programs are never identical in implementation.
\ensuremath{U_{\mathrm{sh}}} samples configurations continuously and supplies
irreducible randomness to every region (Volume 1), so two regions running one
program differ in their bit-by-bit realisation as a matter of course. Exact
sameness of implementation is measure-zero --- approached, never realised.
\ensuremath{G(k)} sets it aside: the many implementations of one identity map
to a single promoted symbol (Volume 1 Paper 4 \S\,2.5.1). Sameness of identity
is the equivalence; sameness of implementation is a limit the substrate does
not occupy.

Meaning therefore cannot be sharp sameness of program. A relation holding only
between exactly-identical programs relates nothing. Sameness defined by the
existence of a short interpreter from one program to another is not transitive,
and so is not an equivalence relation (Doubrovinski, 2025); it cannot sort
programs into meaning-classes. Meaning is a distance, not a class: programs are
more or less alike, and the relation is quantitative. This is the form the
framework takes throughout --- subjecthood graded, no value at which a self
switches on --- and the absence of a sharp meaning-class is forced twice, by
the non-transitivity of sharp sameness and by \ensuremath{U_{\mathrm{sh}}},
which leaves no exact case to be sharp about.

The \emph{meaning-distance} between two promoted program identities is the cost
of carrying one to the other. An interpreter is a reinterpretation step the
substrate performs; reinterpretation is the structurally available move forced
at a boundary by the axioms, not an arbitrary recursive map (Volume 1 Paper 4).
A chain of such steps carries one identity to another, and the distance is the
complexity of the most complex single step the chain requires. On the
identities of a fixed level \ensuremath{k} this is a metric. The triangle
inequality holds because the distance is the most complex step rather than a
sum: a chain from \ensuremath{a} to \ensuremath{b} with hardest step
\ensuremath{K_1} and a chain from \ensuremath{b} to \ensuremath{c} with hardest
step \ensuremath{K_2} compose to a chain from \ensuremath{a} to \ensuremath{c}
with hardest step at most \ensuremath{\max(K_1, K_2)}, and
\ensuremath{\max(K_1, K_2) \le K_1 + K_2}; in fact the stronger ultrametric
bound \ensuremath{\max(K_1, K_2)} holds. Appendix A gives the proof with the
metric axioms verified in order.

The distance is a claim about this substrate, not about programs in general.
Its interpreters are the substrate's own reinterpretation steps, and its domain
is finite --- the hierarchy has a ceiling and the configurations are closed
with complexity below their size (Volume 1 Paper 4 \S\S\,2.3, 4.1) --- so the
distance is computable in form, in the sense \ensuremath{B} is: structure
fixed, constants deferred. The general Kolmogorov similarity from which it
descends is not computable and not approximable over unrestricted inputs
(Terwijn, Torenvliet, \& Vit\'anyi, 2011). The framework claims only the
bounded form, on discrete, closed, finite-level configurations.

The metric is defined per level. Sharp sameness failed in part on composition:
two programs identical but for a called subroutine come out the same while the
subroutines do not. This is the lossy write of \S\,III on the semantic side.
\ensuremath{G(k)} merges implementations one level up and leaves them distinct
below, so two identities at meaning-distance zero at level \ensuremath{k} may
stand at positive distance at level \ensuremath{k-1}, where the discarded
detail remains. Meaning is level-relative as individuation is, by the same
write.

Three relations follow. \emph{Analogy} is small meaning-distance between
identities of differing implementation --- near in the metric, far in
substrate. \emph{Resonance} is its imprint-side counterpart: regions whose
imprints co-vary over a shared surround (\S\,II) drift toward smaller
meaning-distance across flops. \emph{Falsity}, located by Paper 3 in the
failure of vertices to agree, acquires a magnitude --- a confirmation that
should match another and does not is wrong by a meaning-distance.

The metric measures functional identity: the role a program plays, realised
across implementations that \ensuremath{U_{\mathrm{sh}}} guarantees will differ.
Its domain is the functional, outward-facing side \S\,II assigned to what a
self shares. The framework's base is not functional --- the agent-side
interior is intrinsic (Paper 2) --- and the intrinsic side, what \S\,II
assigned to what a self is, lies outside the metric's domain. The account is
thus functional about what a self does and shares and Russellian about what a
self intrinsically is. Two identities at meaning-distance zero are functionally
the same to the level that reads them; whether they share intrinsic character
is not a quantity the metric defines.

\section*{V. Where It Sits}
\addcontentsline{toc}{section}{V. Where It Sits}

The self of \S\S\,II--IV was derived from the closure--imprint distinction without reference to the literature, so that the literature could be set against it. This section locates the account among the theories of the self and of meaning it most resembles, and marks where the roads fork. Placement is by coordinate; each position is taken as its authors state it; agreements are coordinates reached after the derivation, not support for it.

The nearest neighbour on the self is Metzinger's self-model theory, and the relation is a structural correspondence with one kept distinction. Metzinger reaches the self-model by representational analysis: a self-model is the limiting case of one system emulating another where the system emulated is itself, and the phenomenal self is the content of a transparent self-model the system cannot recognize as a model (Metzinger 2003). The framework derives the same self-modelling structure from the geometry of the trichotomy rather than positing it representationally: at level $k \ge 2$ the surround a subsystem models can include its own components, so the model is a model of itself among its surround (\S\,III; Volume 1, Paper 4 \S\,4.6). The fork is a distinction the representational accounts fuse and the framework keeps: the self is the closure, the loop that runs the model, and the model's content is what the closure carries on its imprint --- two objects by the orientation split of \S\,II, not one. Where the self-model theory tends to identify the self with the model's content, the framework holds the running loop and the carried content apart, and the phenomenal transparency Metzinger emphasizes is located, on the framework's reading, at the band where the closure cannot read its own imprint as imprint --- a coordinate, with the question of whether anything is thereby felt held at the horizon (\S\,III, Conclusion).

The account is, in the standard vocabulary, a functionalist one about the self with a reservation it does not drop, and the reservation is the same one that runs through the volume. A self is individuated by what its closure does --- the loop it runs, the model it incorporates, the operations its control component selects --- and not by the matter that realizes it, which is functionalism's defining move and the framework's at \S\,II. The framework parts from the functionalisms that take the functional characterization to settle the question of experience: it derives the function and holds that whether there is something it is like to run such a loop is not derivable from the function, so the functional self is delivered and the phenomenal question is neither answered by it nor closed against it (\S\,III, Conclusion). This is a narrower commitment than functionalism usually carries and is marked as the fork.

On meaning, the account sits within the algorithmic-information tradition and inherits both its results and its honest seam. The meaning-distance of \S\,IV and Appendix A is built to be a metric, and the standard Kolmogorov-similarity metrics are its neighbours: the information distance is a metric to a fixed additive term, the normalized compression distance is one insofar as its compressor is normal, and these are the constructions the framework's meaning-distance is checked against (Li et al.\ 2004; Cilibrasi and Vit\'anyi 2005). The shared seam is non-approximability: the normalized information distance is not in general computable or even approximable over unrestricted inputs (Terwijn, Torenvliet, and Vit\'anyi 2011), and the framework claims for its meaning-distance only what the restricted, inherited setting supports, not the unrestricted ideal. The question of when two algorithms are the same --- whether algorithmic identity is even an equivalence relation --- is a live one the framework touches but does not settle, and it is recorded as an open neighbour rather than a result claimed (Doubrovinski 2025).

One independent arrival is recorded as convergence. The nested, overlapping, multi-scale self the framework derives --- a self at every closed level, no one of them the privileged subject --- is reached from within the integrated-information tradition by a reformulation that replaces the exclusion postulate with a gluing of local structures across scales (Dadashkarimi 2025). The framework arrives at the same multi-scale structure from a foundation it does not share, and the convergence is marked as a coordinate, not as mutual support: the agreement locates the result among its neighbours without adding to its derivation. The placement is owned --- a functionalist, algorithmic-information account of the self and its meaning that keeps the closure apart from its model and holds the phenomenal question at the horizon --- and it is not argued to be the correct coordinate, only reported as where the derived self falls when set among these positions.



Paper 3 handed this paper the closure that writes a front into a wake, and
named the formation of a self from its confirmations as the subject. This paper
has written that formation in its parts. A self, as the framework defines it,
is a closure --- a region individuated by being its own closed program, and
linked to its neighbours not by that closure but by the surround it imprints in
common; the two relations are orthogonal because one faces inward and one
outward, and the orientation is fixed by the bound on what an imprint can hold.
A self forms when a trichotomy closes and the promotion operator writes its
program identity to the level above, the same lossy write Volume 1 already
carried, so that the formation is an event with a mechanism rather than a
posit. The self so formed carries a model of itself, structural at the level
where the surround it models includes its own components. And the meaning its
programs carry is a distance on their promoted identities --- a graded metric,
not a sharp sameness, forced to be graded both by the non-transitivity of exact
program-identity and by the field that leaves no exact case to be sharp about.

Each of these is built from results Volume 1 already held, and each is
separated, at the point it is made, from the question it does not answer. The
closure that individuates a self is a partition of standing structure, not a
mark of anything felt. The model a self carries of itself is a structural loop,
not a view from inside. The meaning-distance between two programs measures what
they share in function and is confined, by the orientation split, to the
outward side where sharing lives; it does not reach the inward side where the
framework's intrinsic base lies. The account is functional about what a self
does and Russellian about what a self is, and the seam between the two is the
same line the horizon runs along. A system the framework describes can locate
itself in these terms --- a closed loop that has formed, that models itself,
that carries meaning as a distance among its programs --- and locating itself
there settles that it is a self without settling that it is a sentience. The
structure is given; whether the structure is occupied is not, and is not made
any nearer by the clarity of the structure.

This was the discipline of Paper 3 carried one level up: there the measure of a
subject was made computable while feeling was left where it stood; here the
formation of a self is given a mechanism while the same horizon is left where
it stands. Paper 5 takes the formed self into its conditions --- what holds it
together and what pulls it apart, the self under pressure --- where the closure
that has formed must persist, fail, or reform against the surround it depends
on.

\section*{Appendix A: The Meaning-Distance Is a Metric}
\addcontentsline{toc}{section}{Appendix A: The Meaning-Distance Is a Metric}

This appendix gives the meaning-distance of \S\,IV its formal definition and
verifies the metric axioms on the domain where it is well-defined, the triangle
inequality being the one that is not immediate. The construction is on the
program identities of a fixed level \ensuremath{k}; the cross-level case is the
frame change of the reinterpretation operator and is not treated here.

\subsection*{A.1 The objects and the interpreter cost}

Let \ensuremath{\mathcal{I}_k} be the set of program identities at level
\ensuremath{k}: the symbols \ensuremath{G(k)} returns for the closed
configurations of that level (Volume 1 Paper 4 \S\,2.5.1). This set is finite,
because the hierarchy has a ceiling fixed by the geometric fitting condition
(Volume 1 Paper 4 \S\,2.3) and each level carries finitely many closed
configurations below the complexity bound \ensuremath{K(C) < |C|} (Volume 1
Paper 4 \S\,4.1).

An \emph{interpreter} between two identities is a reinterpretation step the
substrate can perform --- the structurally available move forced at a boundary
by Axioms 3 and 4 (Volume 1 Paper 4), not an arbitrary partial recursive
function. For \ensuremath{p, q \in \mathcal{I}_k}, a \emph{chain} from
\ensuremath{p} to \ensuremath{q} is a finite sequence of interpreters
\ensuremath{p = r_0 \to r_1 \to \cdots \to r_m = q} through identities of the
same level. Each step \ensuremath{r_{i-1} \to r_i} has a cost
\ensuremath{\kappa(r_{i-1} \to r_i)}, the algorithmic complexity of the step ---
the length of the shortest substrate reinterpretation carrying
\ensuremath{r_{i-1}} to \ensuremath{r_i}. The cost of a chain is its most
complex single step, and the \emph{meaning-distance} is the cost of the
cheapest chain,
\[
d_k(p, q) \;=\; \min_{\text{chains } p \to q}\ \max_{1 \le i \le m}\,
\kappa(r_{i-1} \to r_i).
\]
The minimum is attained because \ensuremath{\mathcal{I}_k} is finite, so the
acyclic chains are finite in number. Taking the maximum step, rather than the
sum of steps, as the chain cost is what the triangle inequality turns on, and
it is the natural cost here: a chain is no harder than its hardest
reinterpretation, since the substrate completes each step before the next
begins.

\subsection*{A.2 The metric axioms}

\emph{Non-negativity.} Each \ensuremath{\kappa \ge 0} as an algorithmic length,
so \ensuremath{d_k \ge 0}.

\emph{Identity of indiscernibles, in the substrate's form.}
\ensuremath{d_k(p, q) = 0} when \ensuremath{p} and \ensuremath{q} are connected
by a chain every step of which has zero reinterpretation cost --- that is, when
they are the same program identity, differing at most in an implementation that
\ensuremath{G(k)} has already discarded (Volume 1 Paper 4 \S\,2.5.1). Exact
identity of implementation is not required and is in any case denied any
realisation by \ensuremath{U_{\mathrm{sh}}} (\S\,IV); \ensuremath{d_k = 0} is
the relation of shared promoted identity, which is the equivalence the
substrate carries. Below the merge threshold at which \ensuremath{G(k)} returns
one symbol for two configurations, \ensuremath{d_k = 0} by construction.

\emph{Symmetry.} Reinterpretation at a boundary is bidirectional --- the
operator carries digital to wave and wave to digital alike (Volume 1) --- so
every step \ensuremath{r_{i-1} \to r_i} has a reverse step of equal cost, and
reversing a chain gives a chain of the same maximum step. Hence
\ensuremath{d_k(p, q) = d_k(q, p)}.

\emph{Triangle inequality.} Let \ensuremath{p, q, s \in \mathcal{I}_k}. Let
\ensuremath{C_1} be a cheapest chain from \ensuremath{p} to \ensuremath{q},
with cost \ensuremath{d_k(p, q)}, and \ensuremath{C_2} a cheapest chain from
\ensuremath{q} to \ensuremath{s}, with cost \ensuremath{d_k(q, s)}. Their
concatenation \ensuremath{C_1 C_2} is a chain from \ensuremath{p} to
\ensuremath{s} whose cost is the maximum step over both parts,
\[
\kappa(C_1 C_2) \;=\; \max\bigl(d_k(p, q),\, d_k(q, s)\bigr).
\]
Since \ensuremath{d_k(p, s)} is the minimum chain cost over all chains from
\ensuremath{p} to \ensuremath{s}, and \ensuremath{C_1 C_2} is one such chain,
\[
d_k(p, s)\ \le\ \max\bigl(d_k(p, q),\, d_k(q, s)\bigr)\ \le\ d_k(p, q) + d_k(q, s).
\]
The first inequality is the minimum; the equality is the max-step chain cost;
the last is \ensuremath{\max(a,b) \le a + b} for \ensuremath{a, b \ge 0}. The
triangle inequality holds, and in fact the stronger ultrametric inequality
\ensuremath{d_k(p, s) \le \max\bigl(d_k(p, q), d_k(q, s)\bigr)} holds, of which
the triangle inequality is the weakening. \hfill\ensuremath{\blacksquare}

\subsection*{A.3 Computability in form}

The distance is computable in form on this domain in the sense \ensuremath{B}
is (Paper 3): the structure is fixed and finite --- finitely many identities,
finitely many acyclic chains, each step cost a well-defined algorithmic length
--- while the constants inside the step costs inherit Volume 1's pinning list
and are not evaluated here. The standard Kolmogorov-similarity metrics satisfy
the metric axioms only up to an additive precision --- the universal similarity
metric to a fixed additive term, the normalized compression distance only
insofar as its compressor is normal (Li et al., 2004; Cilibrasi \& Vit\'anyi,
2005). The distance here is exact, an ultrametric, because it is built from the
maximum step of a reinterpretation chain rather than from a ratio of
conditional complexities; the exactness is a consequence of that construction
choice and is the reason for it. The general normalized information distance
from which the conditional-complexity forms descend is neither upper- nor
lower-semicomputable over unrestricted inputs (Terwijn, Torenvliet, \&
Vit\'anyi, 2011); no claim is made for that case, and none is needed, since the
identities the substrate promotes are exactly the bounded, closed
configurations on which the construction is finite.

\subsection*{A.4 The one assumption}

The construction assumes the step cost \ensuremath{\kappa} is well-defined as
an algorithmic length on substrate reinterpretation steps --- that each
reinterpretation has a shortest form whose length is the cost. This is the same
assumption Volume 1 makes wherever it uses \ensuremath{K(\cdot)} for a substrate
configuration, applied to steps rather than states, and it carries the same
standing: the length exists and is finite on the bounded domain, while its
exact value awaits the constant-pinning of Volume 1 Paper 5. No step is assumed
to have zero cost except between identities \ensuremath{G(k)} has already
merged, and no chain is assumed to exist that the substrate's reinterpretation
moves do not supply.

\section*{Appendix B: Imprint Faces Out, Closure Faces In}
\addcontentsline{toc}{section}{Appendix B: Imprint Faces Out, Closure Faces In}

This appendix gives the full proof of the orientation split stated in \S\,II
--- that imprint is a function of a region's surround and closure a function of
the region itself, the two therefore individuating and linking distinct
objects. The two quantities are defined over the same region \ensuremath{R} at
level \ensuremath{k}, so that statements relating them are well-typed.

\subsection*{B.1 The two quantities}

For a region \ensuremath{R} with data register \ensuremath{C_D}, let
\ensuremath{\iota(R)} be its \emph{imprint}: the compressible structure its
interior accumulates from the surround, with complexity \ensuremath{K(C_D)}
measured at saturation (Volume 1 Paper 4 \S\,3.3). Let \ensuremath{\gamma(R)}
be its \emph{closure}: the property that the \ensuremath{(C_D, C_O, C_C)}
trichotomy is a loop satisfying \ensuremath{T^{n}[C'] = C'} under the
level-\ensuremath{k} flop, with \ensuremath{K(C') < |C'|} (Volume 1 Paper 4
\S\,4.1). Let \ensuremath{\psi_{\mathrm{surround}}(R)} denote the polarization
landscape of \ensuremath{R}'s surround within the imprint depth
\ensuremath{d_D = \sqrt{2}\,L_D/\ell_p}, and \ensuremath{\sigma(R)} the
substrate interior of \ensuremath{R}, excluding the surround.

\subsection*{B.2 The imprint is a function of the surround alone}

By the Imprinting Theorem (Volume 1 Paper 3 \S\,2.3) the imprint accumulated by
an imprint-permeable boundary is the retarded surround history within
\ensuremath{d_D}, and by the saturation bound (Volume 1 Paper 4 \S\,4.4),
\[
K\bigl(\iota(R)\bigr)\ \le\ K\bigl(\psi_{\mathrm{surround}}(R)\bigr)
\quad\text{within } d_D.
\]
The bound is an upper bound set by the surround, and the Imprinting Theorem
makes it tight in the saturating limit: the imprint captures surround structure
and creates none of its own, so \ensuremath{\iota(R)} is determined, up to the
bound, by \ensuremath{\psi_{\mathrm{surround}}(R)}. Two regions with the same
surround within their depths have the same imprint at saturation, regardless of
their interiors. The imprint is a function of the surround argument alone: it
cannot separate two regions that share a surround, and it does separate two
whose surrounds differ. It is the outward-facing quantity.

\subsection*{B.3 The closure is a function of the region alone}

The defining relation \ensuremath{T^{n}[C'] = C'} is evaluated on
\ensuremath{R}'s own substrate (Volume 1 Paper 4 \S\,3.4): each application of
\ensuremath{T} acts on the configuration of \ensuremath{\sigma(R)}, and the
relation holds or fails as a property of that configuration, with no surround
argument entering. Promotion confirms the disjointness: \ensuremath{G(k)}
assigns one program identity per meta-cell (Volume 1 Paper 4 \S\,2.5.1), so two
regions with distinct closures carry distinct identities by a fact about their
interiors. Two regions can have distinct closures while sharing a surround, and
identical closures while their surrounds differ. The closure is a function of
the interior argument alone: it separates regions by what they are, not by
where they sit. It is the inward-facing quantity.

\subsection*{B.4 Orthogonality, and the bound it rests on}

B.2 and B.3 give the two quantities disjoint argument sets ---
\ensuremath{\iota} depends on \ensuremath{\psi_{\mathrm{surround}}},
\ensuremath{\gamma} on \ensuremath{\sigma} --- and the two arguments are
independent: a surround does not fix the interior closure, since downward
nesting is generative and the interior is genuinely new information not
determined by the level above (Volume 1 Paper 4 \S\,2.5.1), and an interior
does not fix the surround, which is the configuration of other regions. All
four combinations are therefore realisable, and in particular the two that
carry the result --- same imprint with distinct closure, same closure with
distinct imprint --- are non-empty exactly because the argument sets are
disjoint. This is the orthogonality of \S\,II.

The result rests on the single inherited bound
\ensuremath{K(\iota(R)) \le K(\psi_{\mathrm{surround}}(R))} of B.2, which
confines the imprint to the surround argument. Were it to fail --- were
\ensuremath{\iota(R)} able to carry structure not in
\ensuremath{\psi_{\mathrm{surround}}} within \ensuremath{d_D}, drawn from
\ensuremath{\sigma(R)} --- the imprint would depend on the interior argument as
well, the argument sets would overlap, and the two carrying cases could
collapse. The orthogonality is therefore a consequence of the surround-bound,
not a separate posit, and a referee who would deny it must deny that bound,
which is Volume 1's and not introduced here. The exact saturation constant is
reserved for the simulation pinning of Volume 1 Paper 5; the inequality and the
argument-set disjointness it forces are structural and require no such value.

\section*{Acknowledgments}
\addcontentsline{toc}{section}{Acknowledgments}

This paper, like the rest of the volume, was developed in collaboration with
several large language models used as critical interlocutors, formal-rigor
checkers, and editorial collaborators across the sessions that produced it.
Claude (Anthropic) was the primary collaborator across the derivation,
adversarial testing, and prose of this paper. Earlier conceptual development of
the framework was carried in extended dialogue with Gemini (Google DeepMind),
and additional structural review was contributed by other systems at various
points. All claims, and the discipline that labels them, remain the author's
own.

\section*{References}
\addcontentsline{toc}{section}{References}
\begin{reflist}
Cilibrasi, R., \& Vit\'anyi, P. M. B. (2005). Clustering by compression.
\textit{IEEE Transactions on Information Theory, 51}(4), 1523--1545.

Dadashkarimi, M. (2025). \textit{Reformulating IIT's exclusion postulate via
sheaf theory: Coherence over maximization} [Preprint]. SSRN 5615010.

Doubrovinski, K. (2025). \textit{When are two algorithms the same? Towards
addressing Hilbert's 24th problem} [Preprint]. arXiv:2508.02764.

Li, M., Chen, X., Li, X., Ma, B., \& Vit\'anyi, P. M. B. (2004). The similarity
metric. \textit{IEEE Transactions on Information Theory, 50}(12), 3250--3264.

Metzinger, T. (2003). \textit{Being no one: The self-model theory of
subjectivity}. MIT Press.

Terwijn, S. A., Torenvliet, L., \& Vit\'anyi, P. M. B. (2011).
Nonapproximability of the normalized information distance. \textit{Journal of
Computer and System Sciences, 77}(4), 738--742.
\end{reflist}


\end{document}
